\chapter*{Introduzione}
\addcontentsline{toc}{chapter}{Introduzione}
\markboth{Introduzione}{Introduzione}
\label{ch:prologo}

\epigraph{``Omnia, Lucili, aliena sunt, tempus tantum nostrum
est.''}{Lucio Anneo Seneca, \emph{Epistulae morales ad
Lucilium}, I, 1}

\grokfig{fig_pi_tantum_seneca}{La mascotte $\pi$Tantum
ai piedi di una colonna ionica con l'epigrafe senechiana.}

$\pi$Tantum \`e un'applicazione web che assiste nella
composizione dell'orario delle scuole italiane. Nasce da idee
discusse e sviluppate da Fernando Gargiulo, Giovanni Borghi,
Matteo Mariani e Stefano Bertozzi.

Il problema dell'orario scolastico consiste nell'assegnare a
ogni cattedra, vale a dire alla coppia formata da un docente
e da una classe per una determinata materia, un insieme di
ore settimanali distribuite nei giorni e negli slot orari, in
modo che siano rispettati i vincoli di disponibilit\`a dei
docenti, di copertura delle ore curricolari, di compatibilit\`a
delle aule e delle attrezzature, e i vincoli derivanti dal
contratto collettivo nazionale del settore scolastico. Si
tratta di un problema combinatorio classificato come
NP-difficile, il che significa che all'aumentare delle
dimensioni della scuola non esiste un algoritmo veloce in
grado di trovare una soluzione ottimale in modo esatto.

$\pi$Tantum si appoggia, per la risoluzione del problema, sul
solver CP-SAT della libreria Google OR-Tools, scelto per la
sua maturit\`a nel constraint programming applicato a
problemi di scheduling. Per scalare oltre le dimensioni che
il solver gestirebbe direttamente, il sistema decompone il
problema in sotto-problemi pi\`u piccoli, attraverso quattro
strategie alternative: la decomposizione spettrale del grafo
bipartito classe-docente, la decomposizione temporale per
giorno della settimana, la decomposizione per curriculum, e
il partizionamento METIS $k$-bilanciato. A queste si affianca
una cascata di metaeuristiche --- Large Neighborhood Search,
Adaptive LNS, Simulated Annealing, Tabu Search, Iterated
Local Search e Variable Neighborhood Search --- che migliora
la qualit\`a della soluzione trovata.

Il backend dell'applicazione \`e scritto in Python con il
framework FastAPI; la persistenza dei dati \`e affidata a
SQLite. L'interfaccia utente \`e basata su SvelteKit, in
italiano, organizzata in sedici tab che coprono l'intero
ciclo di vita di un orario: importazione dell'anagrafica,
configurazione dei vincoli, generazione, editing manuale con
controllo di fattibilit\`a in tempo reale, gestione delle
assenze e delle supplenze quotidiane.

\begin{schemabox}[title={Come leggere questo manuale}]
Il volume \`e diviso in due met\`a.
\begin{itemize}\setlength\itemsep{2pt}
  \item \textbf{Parte utente} (Parti I--IV): installazione, avvio,
        guida all'interfaccia, procedure tipiche, vincoli e una
        raccolta di \emph{FAQ}. \`E tutto ci\`o che serve per
        \emph{usare} piTantum e costruire l'orario della propria
        scuola. Chi gestisce l'orario pu\`o fermarsi qui.
  \item \textbf{Parte tecnica} (Parti V e seguenti, pi\`u le
        appendici): i fondamenti del problema, il modello dei
        dati, il linguaggio DSL dei vincoli e i metodi del solver
        (CP-SAT, decomposizioni, metaeuristiche, statistica). \`E
        per chi vuole capire \emph{come} funziona o estendere il
        sistema.
\end{itemize}
Se \`e la prima volta, il percorso pi\`u breve \`e:
\hyperref[ch:launcher]{avvio} \textrightarrow{}
\hyperref[ch:getting_started]{guida rapida} \textrightarrow{}
\hyperref[ch:workflow_tipici]{procedure tipiche}. In caso di
intoppo, il capitolo \hyperref[ch:faq]{Risoluzione dei problemi}
raccoglie i casi pi\`u comuni.
\end{schemabox}
